% PUT this after day 15 (i.e., do feature selection before assoc analysis)

\documentclass[12pt]{article}

\addtolength{\textwidth}{.8in}
\addtolength{\hoffset}{-.4in}
\addtolength{\textheight}{1in}
\addtolength{\voffset}{-1in}

\usepackage{amsmath}
\usepackage{graphicx}
\usepackage{upquote} % to convert funny quotes to straight quotes
\usepackage{amsfonts}
\usepackage{multicol}
\usepackage{paralist}
\setlength{\parindent}{15pt}

\renewcommand*{\thefootnote}{\fnsymbol{footnote}}


\begin{document}

\begin{center}
\begin{Large}
DATA 219 Notes -- Apr. 12, 2018
\end{Large}
\end{center}

\section{Association Analysis}

Classification isn't the only machine learning task. Another one is called
\textbf{association analysis}, in which the goal isn't to predict the class of
new data items, but rather to determine which features of a data set have
``interesting" relationships. Usually, ``interesting" means ``correlated in a
way that can be taken advantage of."

The principles of Association Analysis are used in fields as diverse as
bioinformatics, recommender systems, and cybersecurity intrusion detection.
Another common application is \textbf{market basket analysis}, in which it can
be employed to detect which sets of products tend to be bought together. This
information is useful, for instance, to retailers who seek optimal product
placement and effective marketing campaigns.

\subsection{Market basket analysis}

All competitive grocery stores today keep track of their customers' purchases
as they check out. Even if they don't have a rewards program to lure customers
into getting a store-specific ``VIP card," they still know which sets of
products were bought together in a single trip to the store. Each such trip,
consisting of multiple products being purchased, is called a
\textbf{transaction}.

This type of association analysis relies on a huge matrix (DataFrame) of 1's
and 0's. Each row is a transaction, and each column is a product the store
sells. A 1 in an entry indicates that \textit{at least one} of that product
was bought in that transaction. (We typically don't store quantities, but just
binary information about ``did this purchase contain that item or not?")
Here's one such data set:

\begin{center}
\begin{tabular}{|c|c|c|c|c|c|}
\hline
diet\_soda & apples & beer & chips\_ahoy & diapers & cat\_food \\
\hline
0 & 1 & 0 & 0 & 0 & 0 \\
1 & 1 & 0 & 0 & 0 & 1 \\
1 & 1 & 1 & 0 & 1 & 0 \\
0 & 0 & 1 & 1 & 1 & 0 \\
1 & 1 & 0 & 1 & 1 & 0 \\
1 & 1 & 1 & 1 & 1 & 1 \\
0 & 0 & 0 & 0 & 0 & 1 \\
. & . & . & . & . & . \\
\hline
\end{tabular}
\end{center}

The first customer in this set just ran into the store to grab a bag of
apples, while the second one bought some Diet Coke, an apple, and a couple
bags of cat food, \textit{etc.}

\vspace{.1in}
Now, some terms:

\begin{itemize}
\itemsep.1em

\item \textbf{itemset}: An itemset is a group of one or more items. For
instance, ``\{apples,beer,cat\_food\}" is an itemset, as is ``\{diet\_soda\}".

\item \textbf{support count}: The support count for a given itemset is
\textit{the fraction of transactions in the entire data set that contain that
itemset.} In the sample data above, the support count for ``\{diet\_soda\}" is
$\frac{4}{7} = .571$, since four out of the seven transactions contained
diet soda. The support for ``\{apples,chips\_ahoy,diapers\}" is just $\frac{2}{7}
= .286$.

\item \textbf{association rule}: An association rule is a statement of a
relationship between two itemsets. It looks like this: ``\{A,B\} $\Rightarrow$
\{C\}," which means ``people who buy A and B tend to also buy C." The itemset
before the arrow is called the \textbf{antecedent} and the one after it is
called the \textbf{consequent}. Association
rules are general correlations, not hard-and-fast mandates. Their reliability
can be measured with the metrics below.

\item \textbf{support}: An association rule's ``support" is the fraction of
transactions in the data set that contain \textit{both} its antecedent
\textit{and} its consequent. It's a measure of how widespread and applicable
it is. The rule ``\{diet\_soda,apples\} $\Rightarrow$ \{beer\}" has a support of
just $\frac{2}{7} = .286$ in the data set above. A rule's support is always
between 0 and 1.

\item \textbf{confidence}: An association rule's ``confidence," on the other
hand, is a measure of its reliability. It's \textit{the fraction of the times
that it could apply where it's actually correct}. In terms of support, it's
$\frac{\textrm{support(whole rule)}}{\textrm{support\_count(antecedent)}}$
The rule ``\{diet\_soda,apples\} $\Rightarrow$ \{beer\}" has a support of
$\frac{2}{4} = .5$ in the data set above, since out of the four times that
diet soda and apples were purchased together, beer was also purchased in two
of them. A rule's confidence is always between 0 and 1.

\item \textbf{lift}: An association rule's ``lift" is a numerical measure of
how much predictive power it has. The key concept is that we want to know how
likely the consequent is to be satisfied when the antecedent is satisfied 
\textit{relative to how likely the consequent is to be satisfied anyway.}
Quantitatively, it's $\frac{\textrm{confidence(whole
rule)}}{\textrm{support(consequent)}}$. A rule's lift is always \textit{between 1 and
$\infty$.} A lift of 1 means it's essentially useless: knowing the antecedent
is true doesn't really make the consequent any more likely.
\end{itemize}

If we crunch all the numbers, we want to look for association rules with high
support, high confidence, and high lift. Which of those three end up being
the greatest priority depends somewhat on the application. If a rule doesn't
have high support, it doesn't come up much, so it may not be super useful. If
it doesn't have high confidence, then the trend it signifies isn't super
reliable, and we need to be aware of that. And if it doesn't have high lift,
then its antecedent doesn't tell us a whole lot that we didn't already know.

\subsection{Implementation: the \texttt{mlxtend} library}

Surprisingly, \texttt{scikit-learn} doesn't seem to have an association
analysis implementation out of the box. (At least, Stephen couldn't find one.)
There is one in the open source \texttt{mlxtend} library, though, by Sebastian
Raschka \& friends.

Let's take it for a spin. First, we'll generate a synthetic data set with a
known association:

\vspace{-.15in}
\begin{center}
\line(1,0){400}
\end{center}
\footnotesize
\begin{Verbatim}[fontsize=\small,samepage=true,frame=none]
import pandas as pd
import numpy as np

eggs = []
bacon = []
bananas = []
pizza = []
coffee = []

NUM_PTS1 = 1000
NUM_PTS2 = 2000

for _ in range(NUM_PTS1):
    eggs.append(1)
    bacon.append(1)
    bananas.append(np.random.choice([1,0],p=[.6,.4]))
    pizza.append(np.random.choice([1,0],p=[.8,.2]))
    coffee.append(np.random.choice([1,0],p=[.3,.7]))

for _ in range(NUM_PTS2):
    eggs.append(np.random.choice([1,0],p=[.7,.3]))
    bacon.append(np.random.choice([1,0],p=[.1,.9]))
    bananas.append(np.random.choice([1,0],p=[.6,.4]))
    pizza.append(np.random.choice([1,0],p=[.8,.2]))
    coffee.append(np.random.choice([1,0],p=[.3,.7]))

transactions = pd.DataFrame({
    'eggs':eggs,
    'bacon':bacon,
    'bananas':bananas,
    'pizza':pizza,
    'coffee':coffee})
\end{Verbatim}
\vspace{-.2in}
\begin{center}
\line(1,0){400}
\end{center}
\normalsize

We have five items in our store, all of which have different priori
probabilities of being purchased (about 60\% of all transactions have bananas,
80\% have pizza, \textit{etc.}) but we've ensured that in addition to the
purely random ones, we've got a solid \textit{thousand} transactions with both
bacon and eggs. This should guarantee that ``\{bacon\} $\Rightarrow$ \{eggs\}"
comes out to be a meaningful association.\footnote{``...and \{eggs\}
$\Rightarrow$ \{bacon\} too, right?" you ask. Well, not necessarily, or at
least not as strong: since 70\% of the random transactions have eggs anyway,
and only 10\% of them have bacon, we're going to get a lot of bacon-less
egg-full transactions, which will dial down the confidence of that
association. But let's wait and see what \texttt{mlxtend} says.}
\pagebreak
Okay, let's analyze. First, just for giggles, let's shuffle the data frame's
rows so they're in a random order (not strictly necessary, but just on
principle):

\begin{Verbatim}[fontsize=\small,samepage=true,frame=none]
transactions = transactions.sample(frac=1)
\end{Verbatim}

and now let \texttt{mlxtend} do its work:
\begin{center}
\line(1,0){400}
\end{center}
\begin{Verbatim}[fontsize=\small,samepage=true,frame=none]
from mlxtend.frequent_patterns import apriori
from mlxtend.frequent_patterns import association_rules

itemsets = apriori(baskets, min_support=0.1, use_colnames=True)
\end{Verbatim}
\vspace{-.3in}
\begin{center}
\line(1,0){400}
\end{center}

This uses the so-called \textbf{apriori algorithm}. It takes place in two
steps: first, generate all itemsets, and compute their ``support," but only
those for a certain minimum support level (in this case, 0.1). The ``but only"
phrase is for simple combinatorics reasons: if you think about it, the number
of possible itemsets is rather large if we have many columns. In fact, if we
have $N$ columns, there are $2^N$ different itemsets!\footnote{Sample numbers:
if we have 100 different items for sale in our store -- certainly not a very
large number at all -- we'll have 1,267,650,600,228,229,401,496,703,205,376
different itemsets to consider. This is far, far more than the number of atoms
in the universe.} The apriori algorithm is smart about not evaluating them
all, because it uses an easy trick: in order for an itemset to have greater
than (say) .1 support, \textit{all of its subsets} must also have at least
that much support. This lets apriori whittle down greatly the number of
itemsets it needs to consider.

\vspace{.1in}
Let's print out the ten itemsets with the greatest support:
\footnotesize
\begin{center}
\line(1,0){400}
\end{center}
\begin{Verbatim}[fontsize=\small,samepage=true,frame=none]
>>> itemsets.sort_values('support',ascending=False).head(10)
   support                itemsets
  0.816000                  [eggs]
  0.801667                 [pizza]
  0.650667           [eggs, pizza]
  0.593333               [bananas]
  0.486000         [bananas, eggs]
  0.475333        [bananas, pizza]
  0.401667                 [bacon]
  0.387667  [bananas, eggs, pizza]
  0.380667           [bacon, eggs]
  0.316000          [bacon, pizza]
\end{Verbatim}
\normalsize
\vspace{.3in}
\pagebreak
Not surprisingly, the one-item itemsets feature prominently in the list, with
eggs on the very top since not only did all of the first 1000 transactions had
eggs, but about 70\% of the others did as well.

\vspace{.1in}
And now, the second step: generate (and compute stats on) all possible
association rules involving those item sets:
\begin{center}
\line(1,0){400}
\end{center}
\small
\begin{Verbatim}[fontsize=\small,samepage=true,frame=none]
>>> rules = association_rules(itemsets)
>>> rules.sort_values('confidence',ascending=False).head(10)

             antecedants consequents  support  confidence      lift 
                 (bacon)      (eggs) 0.375000    0.953390  1.192731 
          (pizza, bacon)      (eggs) 0.302333    0.947753  1.185680 
        (bananas, bacon)      (eggs) 0.228333    0.947441  1.185289 
 (pizza, bananas, bacon)      (eggs) 0.183333    0.940171  1.176194 
         (coffee, bacon)      (eggs) 0.116000    0.938005  1.173485 
         (coffee, bacon)     (pizza) 0.102667    0.830189  1.036872 
          (eggs, coffee)     (pizza) 0.193667    0.812587  1.014889 
                 (bacon)     (pizza) 0.319000    0.811017  1.012927 
        (bananas, bacon)     (pizza) 0.195000    0.809129  1.010569 
           (eggs, bacon)     (pizza) 0.302333    0.806222  1.006939 
\end{Verbatim}
\begin{center}
\line(1,0){400}
\end{center}
\normalsize

As expected, ``\{bacon\} $\Rightarrow$ \{eggs\}" has a very high confidence
(and support, and lift), as do rules of the form ``\{bacon + various other
stuff\} $\Rightarrow$ \{eggs\}." Even in this synthetic case, however, the
lift values aren't all that high numerically, so don't expect them to be.


